% ============================================================== AI USE DISCLOSURE
\newpage
\section*{Generative-AI use disclosure (course AI policy)}
\addcontentsline{toc}{section}{Generative-AI use disclosure}

\noindent\textbf{Tool.} Cursor (Claude) in VS Code, used locally with the assignment folder and conda environment \texttt{PA1\_DL}.

\noindent\textbf{How it was used.} The assistant was used to understand the handout, to ask for section templates, and to sanity-check wording and numbers. Training, evaluation, file layout, and the submitted claims were the student's. Suggested text was kept only after it matched the handout and the saved result files.

\noindent\textbf{Prompts given by the student} (15 representative questions; not a request to complete the assignment):
\begin{enumerate}[leftmargin=*,itemsep=0.6pt]
\item ``Please explain Assignment~2 in your own words: what is Q1 vs Q2, and what must be in the report vs the GitHub repo?''
\item ``Give me a heading template for Question~1 that matches the PDF (architecture, assumptions, task, data, training, robustness). Do not fill the answers.''
\item ``Which river facts are given, and which ones am I allowed to assume? I want history length and the 60-minute horizon labelled as assumptions if the brief does not fix them.''
\item ``Is future rainfall allowed as an input? Is an already announced reservoir release allowed? Please quote the distinction.''
\item ``For Moving MNIST, may I use torchvision's built-in train/test split, or must I split by video index? I am using conda env \texttt{PA1\_DL}.''
\item ``What should Table~1 contain, and how are speed and the symmetry group $G$ defined in the handout? A checklist is enough.''
\item ``How do I count parameters of a ConvLSTM so I can check that Model~A is 480{,}784?''
\item ``What does the optical-flow figure need to show (inputs, arrows, calibration)? A caption template only.''
\item ``Give me a short markdown template for the Q2.5 prediction file. Remind me that it has to be committed before I run the mirror / rotation / speed tests.''
\item ``Same for Q2.11: a prediction template, and remind me it must be committed before I train B, C, and D.''
\item ``In the PDF, what is the difference between models A, B, C, and D? I will choose D myself; I only want the definitions.''
\item ``Write the symmetrised-forecaster formula from the handout so I can check my implementation of $C_T$ and $E_T$.''
\item ``Give me a report outline (Q1, Q2a, Q2b, tables, figures) with empty bullets. I will write the content.''
\item ``Please shorten my Question~1 headings and suggest how to draw the architecture in TikZ. Do not change the design.''
\item ``My GitHub URL is \url{https://github.com/ihuzafaone/DLSTG-PA2}, name Huzaifa, roll 25280009. Put those in the header only. Then list remaining PDF requirements I may have missed.''
\end{enumerate}

\noindent\textbf{Coding help (summary).} Asked for a repo layout, a Table~1 checklist, a parameter-count check for Model~A, a caption template for the optical-flow figure, a Q1 section template, a reminder of what Model~D is allowed to change, and a report outline. Did not ask the assistant to train models or to invent results.

\noindent\textbf{Help on code.} Example snippets and naming suggestions for data loading, ConvLSTM, metrics, training entry points, and symmetrisation (in this repo: \texttt{src/mmdata.py}, \texttt{src/models.py}, \texttt{src/evalkit.py}, \texttt{src/train.py}, \texttt{src/table1.py}, \texttt{src/fig1\_flow.py}, \texttt{q1/river\_prototype.py}). The student adapted those ideas, ran jobs in \texttt{PA1\_DL} on a local NVIDIA RTX~4060 Laptop GPU (8\,GB), and kept Model~D's dilation and the D$+$B run only after validation screening.

\noindent\textbf{My checks on code.} Confirmed Table~1 against the handout; confirmed Model~A (and the screened Model~D) at 480{,}784 parameters; confirmed that every reported test number matches the JSON files in \texttt{results/}.

\noindent\textbf{Question 1.} The assistant drafted empty six-heading templates and TikZ layout hints. History length $L{=}36$ and the six-step (60\,min) horizon are assumptions A1/A2 in the sense that the brief leaves both unspecified; the section keeps those separate from the given river. Future rain is not an input; an already announced release is. The diagram is the TikZ figure; headings were shortened by the student.

\noindent\textbf{Question 2.} The assistant explained required tables/figures and commented on scripts. The student ran Table~1, Figure~1, training, and evaluation, and checked the numbers. Q2.5 was written before the mirror, rotation, and speed tests; Q2.11 was written before training B, C, and D.

\noindent\textbf{Output kept from the assistant.} Clarifications of what the handout was asking, blank section templates, reminders of formulas already given in the brief, and help tightening wording, plus a checklist used on a final pass. Every number reported here was read off the student's own saved result files, never supplied by the assistant.
